% Generated by tools/run_phase22_tables.py
% Preamble requirements: \usepackage{booktabs,graphicx}
% Table 11 is populated only from the verified Phase 23 evidence artifact.

% Table 1: dataset_characteristics
\begin{table*}[t]
\centering
\caption{Characteristics of the normalized smart-meter development dataset.}
\label{tab:dataset_characteristics}
\scriptsize
\resizebox{\textwidth}{!}{%
\begin{tabular}{lll}
\toprule
Characteristic & Value & Interpretation \\
\midrule
Source & https://www.kaggle.com/datasets/ziya07/smart-meter-electricity-consumption-dataset & Third-party Kaggle listing \\
Local file fingerprint & D48DF940FD29444FB9222FC0B59CFDE632E147CCD1D49DC7BED9D5F9DC88F588 & SHA-256 identifies the exact analyzed CSV \\
Rows / variables & 5,000 / 7 & Seven columns including timestamp and supplied label \\
Time span & 2024-01-01 00:00 to 2024-04-14 03:30 & Approximately 104 days \\
Sampling interval & 30 min & 0 non-30-min gaps; 0 duplicate timestamps \\
Missing cells & 0 & No blank cells in the local CSV \\
Supplied labels & Normal 4,750; Abnormal 250 & Abnormal prevalence 5.00\% \\
Rows after maximum lag & 4,664 & 336 half-hour samples represent seven days \\
Measurement units & Not recoverable & Numeric fields are normalized; physical kWh/weather units are not established \\
\bottomrule
\end{tabular}%
}
\vspace{2pt}
\begin{minipage}{0.98\textwidth}
\footnotesize \textit{Note:} The CSV is suitable for a reproducible normalized case study and hardware-in-the-loop portability experiment. It is not sufficient by itself for claims stated in physical energy units. Supplied anomaly labels are publisher-described algorithmic labels rather than independently verified electrical events.
\end{minipage}
\end{table*}

% Table 2: forecasting_models
\begin{table*}[t]
\centering
\caption{Forecasting models and prespecified configurations.}
\label{tab:forecasting_models}
\scriptsize
\resizebox{\textwidth}{!}{%
\begin{tabular}{lll}
\toprule
Model & Configuration & Role \\
\midrule
Compact edge linear & Standardized least squares; 8 features; 2,984 fit + 747 calibration; FP32 deployment & Selected \\
Linear regression & OLS with intercept; 3,731 training rows & Reference \\
Random forest & 200 trees; depth 12; min leaf 2; seed 42 & Benchmark \\
Histogram gradient boosting & lr=0.05; iter=200; leaves=31; min leaf=20; L2=0.1 & Benchmark \\
Seasonal naive (lag 48) & Previous-day value at 30-min sampling (lag 48) & Benchmark \\
Persistence (lag 1) & Previous observation (lag 1) & Benchmark \\
\bottomrule
\end{tabular}%
}
\vspace{2pt}
\begin{minipage}{0.98\textwidth}
\footnotesize \textit{Note:} All learned models use the same leakage-aware lag/calendar inputs in the main comparison. The compact model reserves part of the chronological training window for residual-threshold calibration and is exported as FP32 constants.
\end{minipage}
\end{table*}

% Table 3: chronological_split_configuration
\begin{table*}[t]
\centering
\caption{Chronological holdout and walk-forward split configurations.}
\label{tab:chronological_split_configuration}
\scriptsize
\resizebox{\textwidth}{!}{%
\begin{tabular}{llllll}
\toprule
Protocol & Fold & Train rows & Calibration rows & Test rows & Test interval \\
\midrule
Single holdout & --- & 3,731 & --- & 933 & 2024-03-25 17:30 to 2024-04-14 03:30 \\
Compact final fit/calibration & --- & 2,984 & 747 & 933 & 2024-03-25 17:30 to 2024-04-14 03:30 \\
Expanding walk-forward & 1 & 2,264 & 400 & 400 & 2024-03-03 12:00 to 2024-03-11 19:30 \\
Expanding walk-forward & 2 & 2,664 & 400 & 400 & 2024-03-11 20:00 to 2024-03-20 03:30 \\
Expanding walk-forward & 3 & 3,064 & 400 & 400 & 2024-03-20 04:00 to 2024-03-28 11:30 \\
Expanding walk-forward & 4 & 3,464 & 400 & 400 & 2024-03-28 12:00 to 2024-04-05 19:30 \\
Expanding walk-forward & 5 & 3,864 & 400 & 400 & 2024-04-05 20:00 to 2024-04-14 03:30 \\
\bottomrule
\end{tabular}%
}
\vspace{2pt}
\begin{minipage}{0.98\textwidth}
\footnotesize \textit{Note:} All splits preserve temporal order. Walk-forward test windows are non-overlapping and contain 400 rows each (2,000 total). Each lag uses only an observation available before its prediction timestamp.
\end{minipage}
\end{table*}

% Table 4: forecasting_results
\begin{table*}[t]
\centering
\caption{Single chronological holdout forecasting results.}
\label{tab:forecasting_results}
\scriptsize
\resizebox{\textwidth}{!}{%
\begin{tabular}{lllll}
\toprule
Model & MAE & RMSE & sMAPE (\%) & Observed fit time (s) \\
\midrule
Compact edge linear & 0.130784 & 0.162390 & 39.802 & 0.0114 \\
Linear regression & 0.130960 & 0.162529 & 39.862 & 0.0228 \\
Random forest & 0.132761 & 0.165213 & 40.179 & 1.2757 \\
Histogram gradient boosting & 0.136093 & 0.170102 & 40.953 & 0.8756 \\
Seasonal naive (lag 48) & 0.183463 & 0.230878 & 56.380 & 0.0000 \\
Persistence (lag 1) & 0.188160 & 0.232248 & 57.969 & 0.0000 \\
\bottomrule
\end{tabular}%
}
\vspace{2pt}
\begin{minipage}{0.98\textwidth}
\footnotesize \textit{Note:} MAE and RMSE are in normalized target units. Fit times are observed development-machine timings and are not embedded-device inference times. The compact model uses fewer fitting rows because 747 earlier rows are reserved for anomaly calibration.
\end{minipage}
\end{table*}

% Table 5: walk_forward_results
\begin{table*}[t]
\centering
\caption{Five-fold expanding-window forecasting performance.}
\label{tab:walk_forward_results}
\scriptsize
\resizebox{\textwidth}{!}{%
\begin{tabular}{lllll}
\toprule
Model & MAE mean $\pm$ SD & MAE 95\% CI & RMSE mean $\pm$ SD & Pooled sMAPE (\%) \\
\midrule
Compact edge linear (FP32) & 0.133858 $\pm$ 0.004271 & [0.128555, 0.139161] & 0.165970 $\pm$ 0.004175 & 40.354 \\
Linear regression & 0.133858 $\pm$ 0.004271 & [0.128555, 0.139161] & 0.165970 $\pm$ 0.004175 & 40.354 \\
Random forest & 0.135236 $\pm$ 0.003713 & [0.130626, 0.139846] & 0.167993 $\pm$ 0.003135 & 40.593 \\
Histogram gradient boosting & 0.140116 $\pm$ 0.004235 & [0.134857, 0.145375] & 0.174222 $\pm$ 0.003543 & 41.882 \\
Seasonal naive (lag 48) & 0.188809 $\pm$ 0.008865 & [0.177802, 0.199816] & 0.235346 $\pm$ 0.008103 & 57.707 \\
Persistence (lag 1) & 0.191500 $\pm$ 0.006112 & [0.183910, 0.199089] & 0.237066 $\pm$ 0.007700 & 58.497 \\
\bottomrule
\end{tabular}%
}
\vspace{2pt}
\begin{minipage}{0.98\textwidth}
\footnotesize \textit{Note:} Intervals are two-sided Student-t intervals over five fold-level metrics and should be interpreted descriptively. Metrics use normalized target units. Test windows are non-overlapping; training windows expand.
\end{minipage}
\end{table*}

% Table 6: feature_ablation
\begin{table*}[t]
\centering
\caption{Compact FP32 model feature-group ablation across five temporal folds.}
\label{tab:feature_ablation}
\scriptsize
\resizebox{\textwidth}{!}{%
\begin{tabular}{lllll}
\toprule
Feature group & Features (n) & MAE mean $\pm$ SD & MAE 95\% CI & RMSE mean $\pm$ SD \\
\midrule
Lag + calendar + weather & 11 & 0.133823 $\pm$ 0.004198 & [0.128610, 0.139035] & 0.166092 $\pm$ 0.004093 \\
Weather only & 3 & 0.133848 $\pm$ 0.004323 & [0.128479, 0.139216] & 0.166287 $\pm$ 0.004152 \\
Calendar only & 4 & 0.133852 $\pm$ 0.004395 & [0.128395, 0.139309] & 0.165936 $\pm$ 0.004290 \\
Lag + calendar & 8 & 0.133858 $\pm$ 0.004271 & [0.128555, 0.139161] & 0.165970 $\pm$ 0.004175 \\
Lag + weather & 7 & 0.133865 $\pm$ 0.004203 & [0.128646, 0.139084] & 0.166327 $\pm$ 0.004034 \\
Lag only & 4 & 0.133890 $\pm$ 0.004279 & [0.128577, 0.139203] & 0.166218 $\pm$ 0.004127 \\
\bottomrule
\end{tabular}%
}
\vspace{2pt}
\begin{minipage}{0.98\textwidth}
\footnotesize \textit{Note:} All metrics are in normalized target units. Differences between feature groups are small; the opaque Avg\_Past\_Consumption field is excluded because its construction and leakage properties are undocumented.
\end{minipage}
\end{table*}

% Table 7: anomaly_detection_results
\begin{table*}[t]
\centering
\caption{Residual and isolation-forest anomaly-method agreement with supplied labels.}
\label{tab:anomaly_detection_results}
\scriptsize
\resizebox{\textwidth}{!}{%
\begin{tabular}{lllllllll}
\toprule
Method & TP & FP & FN & TN & Precision & Recall & F1 & Alerts (\%) \\
\midrule
Split conformal (95\%) & 21 & 18 & 31 & 863 & 0.538 & 0.404 & 0.462 & 39 (4.18) \\
Adaptive EWMA residual (p99) & 16 & 14 & 36 & 867 & 0.533 & 0.308 & 0.390 & 30 (3.22) \\
Rolling-week residual (p99) & 10 & 1 & 42 & 880 & 0.909 & 0.192 & 0.317 & 11 (1.18) \\
Static residual (p99) & 8 & 0 & 44 & 881 & 1.000 & 0.154 & 0.267 & 8 (0.86) \\
Isolation forest (5\%) & 19 & 114 & 33 & 767 & 0.143 & 0.365 & 0.205 & 133 (14.26) \\
\bottomrule
\end{tabular}%
}
\vspace{2pt}
\begin{minipage}{0.98\textwidth}
\footnotesize \textit{Note:} These values measure agreement with 52 abnormal and 881 normal publisher-supplied test labels. Those labels are described as Isolation-Forest-generated and are not validated physical anomaly ground truth; therefore the table does not establish real-world event-detection performance.
\end{minipage}
\end{table*}

% Table 8: threshold_sensitivity
\begin{table*}[t]
\centering
\caption{Static residual-threshold sensitivity on the chronological test set.}
\label{tab:threshold_sensitivity}
\scriptsize
\resizebox{\textwidth}{!}{%
\begin{tabular}{llllllllll}
\toprule
Calibration percentile & Threshold & TP & FP & FN & TN & Precision & Recall & F1 & Alert rate (\%) \\
\midrule
90 & 0.284494 & 28 & 46 & 24 & 835 & 0.378 & 0.538 & 0.444 & 7.93 \\
92.5 & 0.310073 & 27 & 31 & 25 & 850 & 0.466 & 0.519 & 0.491 & 6.22 \\
95 & 0.343713 & 21 & 18 & 31 & 863 & 0.538 & 0.404 & 0.462 & 4.18 \\
97.5 & 0.365361 & 17 & 6 & 35 & 875 & 0.739 & 0.327 & 0.453 & 2.47 \\
99 & 0.383026 & 8 & 0 & 44 & 881 & 1.000 & 0.154 & 0.267 & 0.86 \\
99.5 & 0.410987 & 4 & 0 & 48 & 881 & 1.000 & 0.077 & 0.143 & 0.43 \\
\bottomrule
\end{tabular}%
}
\vspace{2pt}
\begin{minipage}{0.98\textwidth}
\footnotesize \textit{Note:} Thresholds are percentiles of absolute residuals on the 747-row chronological calibration set and use normalized units. Classification metrics use the same unverified publisher-supplied labels described in Table 7.
\end{minipage}
\end{table*}

% Table 9: esp8266_numerical_parity
\begin{table*}[t]
\centering
\caption{Numerical and anomaly-decision parity across Python, host C++, and physical ESP8266 execution.}
\label{tab:esp8266_numerical_parity}
\scriptsize
\resizebox{\textwidth}{!}{%
\begin{tabular}{lllllll}
\toprule
Layer & Target & Vectors & Prediction matches & Max. |difference| & Decision matches & Evidence \\
\midrule
Reference & Python float64 & 933 & Reference & Reference & Reference & MAE=0.130784; RMSE=0.162390 \\
Host export & C++17 FP32, x86-64 & 933 & 933/933 & 7.557e-08 & 933/933 & Generated header vs Python reference \\
Physical HIL & ESP8266EX NodeMCU 1.0 & 933 & 933/933 & 8.900e-08 & 933/933 & Replayed held-out vectors; mean kernel time 50.427 $\mu$s \\
\bottomrule
\end{tabular}%
}
\vspace{2pt}
\begin{minipage}{0.98\textwidth}
\footnotesize \textit{Note:} Differences are absolute differences in normalized prediction units. The physical test replays fixed held-out feature vectors stored in program flash; it validates numerical portability and decisions, not live sensor acquisition or field anomaly accuracy.
\end{minipage}
\end{table*}

% Table 10: hardware_resource_utilization
\begin{table*}[t]
\centering
\caption{ESP8266 resource utilization for normal and full-validation firmware.}
\label{tab:hardware_resource_utilization}
\scriptsize
\resizebox{\textwidth}{!}{%
\begin{tabular}{lllll}
\toprule
Resource & Normal used & Normal headroom & Validation used & Validation headroom \\
\midrule
Data RAM & 29,516 / 80,192 (36.81\%) & 50,676 (63.19\%) & 29,852 / 80,192 (37.23\%) & 50,340 (62.77\%) \\
Instruction RAM & 60,343 / 65,536 (92.08\%) & 5,193 (7.92\%) & 60,343 / 65,536 (92.08\%) & 5,193 (7.92\%) \\
Flash code & 257,796 / 1,048,576 (24.59\%) & 790,780 (75.41\%) & 298,132 / 1,048,576 (28.43\%) & 750,444 (71.57\%) \\
Core numeric model state & 104 B & Not a whole-firmware capacity measure & 104 B & 26 FP32 values; excludes code/metadata \\
\bottomrule
\end{tabular}%
}
\vspace{2pt}
\begin{minipage}{0.98\textwidth}
\footnotesize \textit{Note:} Whole-firmware values include the ESP8266 core, Wi-Fi, HTTP, serial reporting, model, and application. Validation vectors are experimental payload rather than model parameters. IRAM is the principal integration constraint, with 5,193 B headroom.
\end{minipage}
\end{table*}

% Table 11: comparison_with_related_studies
\begin{table*}[t]
\centering
\caption{Comparison with closely related forecasting, anomaly-detection, and edge-deployment studies.}
\label{tab:comparison_with_related_studies}
\scriptsize
\resizebox{\textwidth}{!}{%
\begin{tabular}{llllllll}
\toprule
Study & Dataset & Temporal evaluation & Anomaly evidence & Hardware & Full-vector parity & Measured embedded cost & Key limitation \\
\midrule
This work & Normalized 5,000-row case-study dataset & Single holdout + 5-fold expanding window & Agreement with supplied algorithmic labels & Physical ESP8266EX & 933/933 predictions and decisions & 50.427 $\mu$s batch-mean kernel; RAM/IRAM/flash reported & Physical units and real-event ground truth unavailable \\
Shi et al. (2018) [shi2018pooling] & Irish smart-meter data from 920 customers. & Study-specific train/test forecasting; no rolling-origin result used in our comparison & None & TensorFlow development platform; no constrained-device deployment reported. & Not reported & No MCU cost reported & Accuracy study rather than embedded portability or anomaly detection. \\
Kong et al. (2019) [kong2019lstm] & Public real residential smart-meter dataset. & Study-specific residential forecasting experiment; no rolling-origin result used here & None & No constrained-microcontroller execution reported. & Not reported & No MCU cost reported & Focuses forecasting accuracy rather than reproducible software-to-device transfer. \\
Liu and Nielsen (2018) [liu2018scalable] & A real-world dataset and a large synthetic dataset. & Online/iterative streaming evaluation & Real and synthetic pattern-anomaly evaluation & Big-data system architecture; no microcontroller deployment. & Not applicable & Scalability; no MCU RAM/latency & System-level online scalability is not evidence of TinyML feasibility or numerical portability. \\
Utomo and Hsiung (2020) [utomo2020multitiered] & One year of 30-min data from 200 US households; labels based on a statistical rule. & Sliding windows and daily look-ahead labels & Mean + 3 SD statistical labels, expanded over look-ahead windows & Raspberry Pi edge execution. & Not reported & Raspberry Pi latency and model file size & Raspberry Pi is substantially less constrained than ESP8266, and anomaly labels are constructed statistically. \\
Hu and Tang (2020) [hu2020edge] & Numerical smart-meter system experiments. & Offline/online cloud-edge numerical evaluation & Not the primary evaluation & Edge-enabled architecture; a specific ESP-class MCU parity experiment is not reported. & Not reported & Execution-time results; no ESP8266 accounting & Architecture-level evidence does not isolate deterministic microcontroller inference portability. \\
Gajendran et al. (2026) [gajendran2026ihems] & PZEM-004T measurements from two residential halls, sampled every 5 s and aggregated hourly. & Chronological 70/15/15; next-hour prediction & None & ESP8266 sensing/Wi-Fi gateways; forecasting evaluated on i7 CPU and RTX 3060 GPU. & Not reported & PC ML latency/memory; ESP8266 link latency 72--95 ms & The ESP8266 transports measurements but does not execute the forecasting models reported in the performance table. \\
Banbury et al. (2021) [banbury2021mlperf] & Keyword spotting, visual wake words, image classification, and machine-anomaly benchmarks. & Task-specific benchmark splits; not load forecasting & Machine anomaly benchmark, not smart-meter events & Multiple ultra-low-power MCU-class systems. & Not its research question & Latency and energy benchmark framework & No smart-meter load forecasting task or Python-to-export parity question. \\
Hernández et al. (2024) [hernandez2024daily] & Seven months of hourly appliance data from a commercial Wibeee meter in one four-tenant household. & Seven-month household series; next-hour prediction & Activity-alarm labels from an instrumented household & Commercial physical meter collected data; edge implementation is described as a future objective. & Not reported & Edge feasibility discussed; MCU cost not reported & Single household and application-specific activity alarms limit generalization. \\
Li et al. (2024) [li2024federated] & Building Data Genome 2 and Irish CER customer-behaviour data. & First year train; subsequent half-year test; five runs with 95\% CIs & None & Thirty ARM Cortex-M4 MCUs plus three edge PCs and one cloud server. & Not reported as a Python/C++ parity audit & MCU memory, training time, communication & One representative meter hardware configuration and no anomaly-detection study. \\
Kolosov et al. (2025) [kolosov2025instrumental] & Measured 15-kHz voltage/current/power signatures and NILM evaluation data. & NILM train/test scenarios; not load forecasting & None; appliance disaggregation task & Custom analog front end plus six Raspberry Pi/accelerator/FPGA-class platforms. & Cross-platform accuracy, not row-wise export parity & Latency, power, throughput, efficiency across six boards & Targets NILM on substantially stronger hardware, not low-rate forecasting/anomaly decisions on ESP8266. \\
\bottomrule
\end{tabular}%
}
\vspace{2pt}
\begin{minipage}{0.98\textwidth}
\footnotesize \textit{Note:} External rows come from the targeted Phase 23 peer-reviewed gap analysis. Citation keys map to references/references.bib. The search is not a systematic review and does not support a 'first' claim. Metrics are not compared numerically across incompatible datasets/tasks.
\end{minipage}
\end{table*}
